\chapter{Clause structure}

\section{Basic word order}

\subsection{Introduction}

Niklorsk is an analytic language with a basic word order of Subject-Verb-Object. It is a mostly head-initial language, with prepositions, prenominal adjectives, and preverbal adverbs.

In a neutral sentence, the object will always follow the verb. In the case of a ditransitive verb, i.e., a verb that takes both a direct and an indirect object, the indirect object will precede the direct object. For example,

\pex
\begingl
\gla Ag vædar Martin. Ag brar afte Martin.//
\glb \First{}\Sg{} see.\Pst{} Martin. \First{}\Sg{} give.\Pst{} book Martin.//
\glft `I saw Martin. I gave Martin a book.'//
\endgl
\xe

Prepositional phrases generally appear after the subject the verb but before the object. For example,

\pex
\begingl
\gla Ag vædar Martin dänn kinnotek.//
\glb \First{}\Sg{} see.\Pst{} Martin in cinema//
\glft `I saw Martin at the movies.'//
\endgl
\xe

However, if the object is a pronoun, it will precede the prepositional phrase.

\pex
\begingl
\gla Martin vædar han dänn kung.//
\glb Martin see \Refl{} in mirror//
\glft `Martin saw himself in the mirror.'//
\endgl
\xe

\subsection{Constinuent fronting}

Niklorsk allows the movement of some constituents to the beginning of the sentence. THe movement of a constituent to the beginning of the sentence is sometimes used to emphasize the constituent, although in some cases the movement does not carry any particular emphasis and usually has the same currency as the neutral word order.

\pex
\begingl
\gla Matti brar skämd om hond.//
\glb Matti give.\Pst{} toy the.\Dat{} child//
\glft `Matti gave the child a toy.'//
\endgl
\xe
\pex
\begingl
\gla Skämd brar Matti om hond.//
\glb toy give.\Pst{} Matti the.\Dat{} child//
\glft `Matti gave the child a toy.'//
\endgl
\xe
\pex
\begingl
\gla Om hond brar Matti skämd.// 
\glb the.\Dat{} child give.\Pst{} Matti toy//
\glft `Matti gave the child a toy.'//
\endgl
\xe

Personal and reflexive pronouns, however, cannot be fronted in this manner. Thus the sentences in the second example below are ungrammatical.

\pex
\begingl
\gla Matti brar när skämd. Matti brar när ja.//
\glb Matti give.\Pst{} \Third{}\Sg{}.\Dat{} toy Matti give.\Pst{} \Third{}\Sg{}.\Dat{} it//
\glft `Matti gave him a toy. Matti gave it to him.'//
\endgl
\xe

\pex
\begingl
\gla *När brar Matti skämd. *Ja brar Matti när.//
\glb \Third{}\Sg{}.\Dat{} give.\Pst{} Matti toy it give.\Pst{} Matti \Third{}\Sg{}.\Dat{}//
\glft Intended: `Matti gave him a toy. Matti gave it to him.'//
\endgl
\xe


Adverbs of time and manner are usually placed at the beginning of the sentence. This fronting results in the inversion of the usual subject-verb order, with the verb now preceding the subject. The construction is not necessarily emphatic and cooccurs regularly with the normal word order. For example,

\pex
\begingl
\gla Ag vædar Martin i kvett dänn kampus.//
\glb \First{}\Sg{} see.\Pst{} Martin in yesterday the campus//
\glft `I saw Martin on campus yesterday.'//
\endgl
\xe
\pex
\begingl
\gla I kvett vædar ag Martin dänn kampus.//
\glb in yesterday see.\Pst{} \First{}\Sg{} Martin the campus//
\glft `Yesterday I saw Martin on campus.'//
\endgl
\xe

\pex
\begingl
\gla Eg störs hvengaren rontæg par veerser.//
\glb the error find.\Pst{}.\Pass{} quickly by:the scientist//
\glft `The error was quickly found by the scientist.'//
\endgl
\xe

\pex
\begingl
\gla Rontæg hvengaren eg störs par veerser.//
\glb quickly find.\Pst{}.\Pass{} the error by:the scientist//
\glft `The error was quickly found by the scientist.'//
\endgl
\xe

The phrase \nikp{dänn kampus}{on campus} may also be placed at the beginning of the sentence, but such a construction is significantly more marked and would only be used if the emphasis is indeed on the location of the event. Two or more adverbial phrases cannot be fronted at the same time. Thus the following sentence is ungrammatical:

\pex
\begingl
\gla *I kvett dänn kampus vædar ag Martin.//
\glb in yesterday the campus see.\Pst{} \First{}\Sg{}//
\glft Intended: `Yesterday I saw Martin on campus.'//
\endgl
\xe

In general, adverbs expressing location are not fronted except for emphasis. Exceptions to this rules are the adverbs \nikp{men}{here}, \nikp{imår}{there (medial)} and \nikp{imed}{there (distal)} and indefinite pronouns like \nikp{elge käg}{anywhere}, etc., used adverbially, which are usually fronted without changing the sentence's emphasis. Nevertheless, if the aforementioned adverbs occurs in a sentence with an adverb of time or manner, the latter will take precedence in fronting.

\pex
\begingl
\gla Men slötar Martin til hjekks av Kari.//
\glb here come.\Pst{} Martin for speak.\Inf{} with Kari//
\glft `Martin came here to speak with Kari.'//
\endgl
\xe

\pex
\begingl
\gla I kvett slötar Martin men til hjekks av Kari.//
\glb in yesterday come.\Pst{} Martin here for speak.\Inf{} with Kari//
\glft `Yesterday Martin came here to speak with Kari.'//
\endgl
\xe

\pex
\begingl
\gla Elge käg igom ga ferms.//
\glb any where can.\Pres{} \Second.\Sg{} sit.\Inf{}//
\glft `You can sit anywhere.'//
\endgl
\xe

In the case of a modal verb, adverb fronting affects only the modal verb, with the main verb following the subject. For example,

\pex
\begingl
\gla Ag ignar hjekkuts av Martin i kvett dänn kampus.//
\glb \First{}\Sg{} can.\Pst{} speak.\Inf{} with Martin in yesterday the campus//
\glft `I was able to speak with Martin on campus yesterday.'//
\endgl
\xe

\pex
\begingl
\gla I kvett ignar ag hjekkuts av Martin dänn kampus.//
\glb in yesterday can.\Pst{} \First{}\Sg{} speak.\Inf{} with Martin the campus//
\glft `Yesterday I was able to speak with Martin on campus.'//
\endgl
\xe


\section{Questions}

\subsection{Yes-no questions}

Yes-no questions are formed by changing the intonation of the sentence to indicate a question. In writing, a question mark is used at the end of the sentence. Alternatively, the question particle \nk{inte} may be placed before the verb. For example,

\pex
\begingl
\gla Matti inte bördet någ?//
\glb Matti \Q{} eat.\Perf{} already//
\glft `Has Matti already eaten?'//
\endgl
\xe

\nk{Inte} is required in quoted questions, as in the following example. It may be placed either before the main verb or before the verb in the subordinate clause.

\pex
\begingl
\gla Meg inte hvæd, efte Matti bördet någ?//
\glb \Second{}\Sg{} \Q{} know if Matti eat.\Perf{} already//
\glft `Do you know if Matti has already eaten?'//
\endgl
\xe

A single noun/pronoun or noun/pronoun phrase, or the interrogative pronoun by itself, may be used to form a question. The question particle \nk{inte} is required in such cases.

\pex
\begingl
\gla Käll vaedar ga i Malmö? Martin inte?//
\glb who see.\Pst{} \First{}\Pl{} in Malmö Martin \Q{}//
\glft `Whom did we see in Malmö? Was it Martin?'//
\endgl
\xe


\subsection{Content questions}

Content questions (or wh-questions) are formed by placing an interrogative pronoun at the beginning of the sentence (see \S\ref{sec:interrogative-pronouns}). The interrogative pronoun replaces the element in the sentence that is being questioned. Interrogative sentences generally follow the same word order as declarative sentences. However, if the interrogative pronoun replaces the object of the verb (or some other element of the predicate), the verb will precede the subject. Moreover, as with yes-no questions, the question particle \nk{inte} may be placed directly after the interrogative pronoun for emphatic or stylistic reasons. For example,

\pex
\begingl
\gla Käll vædar ægn?//
\glb who see.\Pst{} me//
\glft `Who saw me?'//
\endgl
\xe

\pex
\begingl
\gla Käll (inte) vædar ag?//
\glb who \Q{} see.\Pst{} \First{}\Sg{}//
\glft `Whom did I see?'//
\endgl
\xe

The interrogative pronoun may also be used attributively to further qualify the answer the speaker is expecting, similar to the use of `Which?' in English. For example, if in the following example the speaker is expecting the answer to belong to a specific group of people, viz., the listener's friends, the question would be:

\pex
\begingl
\gla Käll maug (inte) vædar ag?//
\glb who.\Pl{} friend \Q{} see.\Pst{} \First{}\Sg{}//
\glft `Which of (your) friends did I see?' or `Which friend of yours did I see?'//
\endgl
\xe

Note that the plural interrogative pronouns \nikp{käller}{who} and \nikp{känner}{what} are used when the speaker is expecting a plural answer. Thus, if the response is a singular noun, the answer is usually qualified by the adverb \nikp{ngå}.

\pex
\begingl
\gla Käller vædar ægn? --- Martint maugen. --- Martin ngå.//
\glb who see.\Pst{} me --- Martin.\Gen{} friend.\Pl{} --- Martin only//
\glft `Who saw me?' 'Martin's friends.' 'Only Martin.'//
\endgl
\xe




\section{Subordinate clauses}

\subsection{Relative clauses}

A relative clause is a subordinate clause that modifies a noun. A relative clause is introduced by the corresponding interrogative pronoun (cf. \S\ref{sec:interrogative-pronouns}) which functions as a relative pronoun in this context. The verb in a relative clause is placed at the beginning of the clause, directly after the relative pronoun. The word order of the relative clause is otherwise the same as that of a main clause. For example,


\pex
\begingl
\gla Ag vædar en bliv dänn magasin.//
\glb \First{}\Sg{} see.\Pst{} the man in store//
\glft `I saw the man at the store.'//
\endgl
\xe
\pex
\begingl
\gla En bliv, käll vædar ag dänn magasin, ö Martint dur.//
\glb the man who see.\Pst{} \First{}\Sg{} in store, is Martin.\Gen{} boy//
\glft `The man whom I saw at the store is Martin's son.'//
\endgl
\xe
\pex
\begingl
\gla Meg inte jagt en forst, kamm vædar ag Martin dänn magasin?//
\glb you \Q{} remember the.\M{} time when see.\Pst{} \First{}\Sg{} Martin in store//
\glft `Do you remember the time when I saw Martin at the store?'//
\endgl
\xe

\subsection{Complement clauses}
A complement clause is a subordinate clause that functions as the complement of a verb or noun. In Niklorsk, complement clauses are introduced by the complementizer \nk{megd}. The word order of the complement clause is reversed, with the verb appearing at the beginning of the clause directly after the word \nk{megd}.

\pex
\begingl
\gla Martin malmar at Oslo i fänt hlåd.//
\glb Martin go.\Pst{} to Oslo in last summer//
\glft `Martin went to Oslo last summer.'//
\endgl
\xe
\pex
\begingl
\gla Ag sædar, megd malmar Martin at Oslo i fänt hlåd.//
\glb \First{}\Sg{} hear.\Pst{} \Compl{} go.\Pst{} Martin to Oslo in last summer//
\glft `I heard that Martin went to Oslo last summer.'//
\endgl
\xe
